\begin{center}
\begin{minipage}{0.92\textwidth}
\small
\textbf{Аннотация.}
Программное обеспечение специализированных вычислительных ускорителей отлаживают
на функциональных симуляторах, которые должны численно совпадать с устройством.
Самую массовую и медленную операцию~--- матричное умножение~--- в таком
симуляторе выгодно переносить с центрального процессора на графический, однако
из-за особенностей арифметики с плавающей точкой разные реализации дают слегка
различные результаты, и эти расхождения накапливаются по слоям сети. В работе
предложены аналитическая модель накопления расхождений и измерительный конвейер,
раздельно оценивающий вклад компилятора нейросетей и вклад реализации матричного
умножения. Установлено, что главным источником расхождений служит компилятор, а
не выбор реализации умножения, поэтому перенос операции на графический процессор
не ухудшает вычислительную точность относительно эталона на центральном
процессоре. Показано также, что характер накопления определяется типом ошибки
(случайная или систематически смещённая) сильнее, чем её величиной.

\smallskip
\textbf{Ключевые слова:} функциональный симулятор, матричное умножение, GEMM,
вычислительная точность, арифметика с плавающей точкой, компилятор XLA,
нейронные сети, накопление погрешностей.
\end{minipage}
\end{center}

\vspace{0.5em}
